\section{Representing an arbitrary automorphism}
\label{sec:representation}

We invoke the following established result. Its
scope is essential: it applies to any automorphism of the full
degree order and supplies a single function valid at every input.

\begin{theorem}[Slaman--Woodin]\label{thm:representation}
For every $\pi\in\Aut(\DT,\leT)$ there is an arithmetic function
$\widehat F:2^\N\to2^\N$ such that
\[
\degT{\widehat F(A)}=\pi(\degT A)
\qquad\text{for every }A\subseteq\N.
\]
\end{theorem}

This is Theorem~6.3.1(2) of \cite{SlamanWoodin2005};
see also Theorem~4.29 in the author version of
\cite{SlamanGlobal}. The generic-input assertion in the first
clause of the cited Theorem~6.3.1 is separate from the all-input assertion
used here.

Arithmetic here describes the graph of a single total function.
That graph is Borel. By the Borel-graph
criterion for functions between Polish spaces
\cite[Corollary~4.15]{MarkerDST}, $\widehat F$ is Borel.
The weaker Borel conclusion also appears directly as
Corollary~4.25 in the author version of \cite{SlamanGlobal}.

\begin{lemma}\label{lem:lift}
For an arbitrary $\pi\in\Aut(\DT,\leT)$ there is a Borel
function $F:\R\to2^\N$ with countable fibers such that
\begin{align}
\degT{F(x)}&=\pi(d(x))\qquad(x\in\R),\label{eq:representation}\\
F(x+y)&\leT F(x)\oplus F(y)\qquad(x,y\in\R).
\label{eq:addition}
\end{align}
\end{lemma}

\begin{proof}
Take $F=\widehat F\circ C$ from
Theorem~\ref{thm:representation} and Lemma~\ref{lem:cuts}.
If $F(x)=F(y)$, then $\pi(d(x))=\pi(d(y))$, hence $d(x)=d(y)$.
Each fiber of $F$ therefore injects, by $C$, into a single
Turing degree. Such a degree is countable: its members are
among the outputs of countably many oracle programs with
any fixed representative as oracle.

By Lemma~\ref{lem:cuts}, $d(x+y)\leT d(x)\vee d(y)$.
Since $\pi$ preserves order and joins,
\[
\degT{F(x+y)}
\leT \pi(d(x))\vee\pi(d(y))
=\degT{F(x)\oplus F(y)}.
\]
This is exactly \eqref{eq:addition}.
\end{proof}

The representing function is not required to be injective,
or to give the same output set on equivalent input sets.
Equation~\eqref{eq:representation} prescribes only its output
degree. The countable-fiber property is the consequence of
injectivity of the automorphism that will enter the argument.
